\documentclass[twocolumn,pra,amsmath,amssymb,superscriptaddress]{revtex4-2}

\usepackage{graphicx}
\usepackage{hyperref}
\usepackage{booktabs}
\usepackage{multirow}
\usepackage{xcolor}
\usepackage{listings}

\lstset{
  language=Python,
  basicstyle=\ttfamily\footnotesize,
  keywordstyle=\color{blue},
  stringstyle=\color{red!60!black},
  commentstyle=\color{green!50!black},
  breaklines=true,
  frame=single,
  xleftmargin=1em,
  xrightmargin=0.5em,
}

\begin{document}

\title{Structurally Sparse Molecular Hamiltonians for Near-Term Quantum Simulation}

\author{J.~Loutey}
\affiliation{Independent Researcher, Kent, Washington}

\date{\today}

\begin{abstract}
We present a library of qubit Hamiltonians for molecular quantum simulation
that are structurally sparser than conventional Gaussian-basis alternatives
by one to three orders of magnitude. The Hamiltonians are constructed from
a spectral graph theory framework (GeoVac) that exploits natural coordinate
separations in molecular electronic structure, producing block-diagonal
electron repulsion integrals and Gaunt-enforced selection rules that persist
through Jordan-Wigner encoding. For the composed architecture (PK-partitioned),
LiH requires 334 Pauli terms at 30 qubits versus 907 for Gaussian STO-3G
at 12 qubits, with 13$\times$ fewer qubitwise-commuting measurement groups.
The balanced coupled architecture (PK-free) binds LiH with an FCI minimum
at a computed $R_{\mathrm{eq}} = 3.227$~bohr (7.0\% above the experimental
3.015~bohr) at $n_{\max}=2$ (878 Pauli terms, 30 qubits); a single-point
evaluation at the experimental geometry reaches 0.20\% energy error at
$n_{\max}=3$ (84 qubits) when solved exactly.  Chemistry-accuracy scope at $n_{\max}=2$ is
\emph{row-conditional}: first-row hydrides (LiH, BeH$_2$, H$_2$O, $\ldots$)
bind with documented potential-energy-surface (PES) topology, while
second-row hydrides (NaH and below) exhibit monotone overattraction under
every tested configuration --- frozen-core balanced, multi-zeta valence,
screened cross-center, explicit-core Hartree--Fock --- with no interior PES
minimum at the qubit-Hamiltonian level. We document this scope boundary
empirically in \S\ref{subsec:scope_boundary}; the framework's
primary value proposition is structural sparsity with quantitative
truncation-error bounds (Paper~38 state-space GH convergence), not energy-accuracy
parity with CCSD(T) on second-row systems.  The \texttt{hamiltonian()}
registry covers 37 systems (35 composed molecules across three rows of
the periodic table, $Z=1$--36, plus the He and H$_2$ reference atoms),
including multi-center systems
(CO, N$_2$, F$_2$), mixed
frozen-core molecules (NaCl), and the full first transition series as hydrides (ScH through ZnH),
with Pauli scaling $O(Q^{2.5})$
confirmed universal across all three rows, compared to
$O(Q^{3.9\text{--}4.3})$ for Gaussian baselines.  The composed Pauli
count is exactly linear in qubit count (coefficient $11.10$, exact,
for main-group molecules; $9.23$ for $d$-orbital blocks under the
pair-diagonal ERI approximation (Sec.~\ref{sec:eri_rule})---the exact
global-$M_L$ rule makes the $d$-block coefficient higher, not lower),
so adding atoms to a molecule increases quantum cost linearly.
Second- and third-row molecules use analytical frozen cores
([Ne], [Ar], [Ar]3$d^{10}$); the frozen-core energy enters only
through the identity operator and does not affect quantum measurement
cost.  Isostructural molecules with frozen cores produce identical
Pauli counts across periodic table rows.  The Hamiltonians are available
as a Python library with Qiskit, OpenFermion, and PennyLane export.
\end{abstract}

\maketitle

%% ============================================================
\section{Introduction}
\label{sec:intro}

Quantum simulation of molecular electronic structure is among the most
promising near-term applications of quantum computing~\cite{McArdle2020,
Bauer2020,Cao2019}. However, the cost of implementing molecular Hamiltonians
on quantum hardware is dominated by the number of Pauli terms in the
Jordan-Wigner (or equivalent) encoding: each distinct Pauli string requires
separate measurement circuits, and the total number of shots scales with
the square of the Pauli 1-norm~\cite{Wecker2015,Rubin2018}.

For a molecule encoded in $Q$ qubits using a Gaussian atomic orbital basis
and the Jordan-Wigner transformation, the Pauli term count scales as
$O(Q^4)$ due to the dense four-index electron repulsion integrals (ERIs).
Post-hoc compression techniques---double factorization~\cite{vonBurg2021},
tensor hypercontraction~\cite{Lee2021}, qubit tapering~\cite{Bravyi2017}---can
reduce this cost but do not alter the underlying structure of the Hamiltonian.

An alternative is to construct Hamiltonians that are \emph{structurally}
sparse: the ERIs are block-diagonal by construction, with zeros enforced
by angular momentum selection rules rather than discovered by compression.
The GeoVac framework achieves this by exploiting the natural coordinate
separations of molecular electronic structure~\cite{GeoVac_Paper7,
GeoVac_Paper14,GeoVac_Paper17,GeoVac_Paper19}. Each electron group
(core, bond pair, lone pair) is described in its own natural coordinate
system, coupled via cross-center nuclear attraction integrals and
cross-block Coulomb integrals.  The angular matrix elements are computed
from Gaunt integrals (Wigner 3$j$ symbols), which impose exact selection
rules that translate directly into structural zeros in the qubit Hamiltonian.

The result is Pauli scaling of $O(Q^{2.2\text{--}2.5})$ across first-row
molecules, with 51--1{,}712$\times$ fewer terms than Gaussian raw
Jordan-Wigner baselines at comparable qubit counts~\cite{GeoVac_Paper14}.
This structural sparsity is compatible with all downstream optimizations
(tapering, grouping, tensor factorization) and is basis-intrinsic rather
than compression-dependent.

The relationship to numerical factorization techniques can be sharpened.
Double factorization~\cite{vonBurg2021} applied to a GeoVac ERI tensor
reproduces the analytic multipole-channel decomposition exactly at the
production basis ($n_{\max} = 2$): the double-factorization rank equals
the count of distinct radial-density channels (seven for valence
$\{1s, 2s, 2p\}$), and every singular vector of the reshaped two-electron
tensor lives strictly inside the multipole-channel subspace.  At larger
basis ($n_{\max} \geq 3$), double factorization still produces only
linear combinations of multipole channels, but the radial-integral matrix
develops a graded singular-value spectrum reflecting the well-known
near-degeneracy of hydrogenic radial integrals; the additional
compression beyond the bare multipole count is therefore a numerical
realization of structure already present in the analytic decomposition,
not a new algebraic substructure.  Double factorization on a GeoVac
Hamiltonian is, in this precise sense, the multipole expansion accessed
by SVD rather than by closed form (see Paper~14 \S\ref{sec:intro} and
\texttt{debug/sprint\_df\_multipole\_lift\_memo.md}).

This paper presents benchmark data for the GeoVac Hamiltonian library:
accuracy characterization via full configuration interaction (FCI),
quantum resource comparison against Gaussian baselines, and installation
instructions.  Theory derivations are deferred to the supporting
papers~\cite{GeoVac_Paper7,GeoVac_Paper14,GeoVac_Paper17,GeoVac_Paper19}.

%% ============================================================
\section{The Hamiltonian}
\label{sec:hamiltonian}

\subsection{What the user receives}

A GeoVac Hamiltonian is a standard qubit operator---a weighted sum of
Pauli strings---delivered as a \texttt{SparsePauliOp} (Qiskit),
\texttt{QubitOperator} (OpenFermion), or \texttt{qml.Hamiltonian}
(PennyLane):

\begin{lstlisting}
from geovac.ecosystem_export import hamiltonian
h = hamiltonian('LiH', R=3.015, max_n=2)
op = h.to_qiskit()    # 334 Pauli terms
op_of = h.to_openfermion()
op_pl = h.to_pennylane()
print(h.n_qubits)     # 30
print(h.one_norm)      # 33.3 Ha
\end{lstlisting}

The operator is ready for use in any VQE, ADAPT-VQE, QPE, or other
quantum algorithm that consumes a Pauli Hamiltonian.

\subsection{Two architectures}

GeoVac provides two Hamiltonian architectures for each molecule:

\paragraph{Composed (PK-partitioned).}
The Phillips-Kleinman (PK) pseudopotential, a one-body operator that
enforces core-valence orthogonality, is separated from the qubit
Hamiltonian and evaluated classically from the VQE one-particle
reduced density matrix (1-RDM)~\cite{GeoVac_Paper17}:
\begin{equation}
  E_{\text{total}} = \langle\psi|H_{\text{elec}}|\psi\rangle
    + \mathrm{Tr}(h^{\text{PK}}\,\gamma)\,.
\end{equation}
This produces the sparsest qubit operator (334 Pauli terms for LiH)
at the cost of requiring a classical post-processing step.

\paragraph{Balanced coupled (PK-free).}
Cross-center nuclear attraction integrals, computed via a multipole
expansion that terminates exactly at $L = 2\ell_{\max}$ by Gaunt
selection rules~\cite{GeoVac_Paper19}, replace the PK pseudopotential.
This produces a larger operator (878 Pauli terms for LiH at
$n_{\max}=2$) but one that is variationally complete within the
composed partition---FCI on this Hamiltonian gives physical energies
without any classical correction.

\subsection{How it works (brief)}

Electrons are partitioned into groups (core, bond pairs, lone pairs)
based on the chemical structure.  Each group is described in its natural
coordinate system (hyperspherical for core electrons, molecule-frame
hyperspherical for bond pairs).  Within each group, the angular basis
states are labeled by hydrogenic quantum numbers $(n, \ell, m)$, and
all angular matrix elements are evaluated from Gaunt integrals---products
of spherical harmonics integrated analytically via Wigner 3$j$ symbols.
The selection rules $|\Delta\ell| \leq k$, $|\Delta m| \leq k$ produce
exact structural zeros in the ERI tensor that persist through
Jordan-Wigner encoding.  Full derivations are in
Refs.~\cite{GeoVac_Paper7,GeoVac_Paper14,GeoVac_Paper17,GeoVac_Paper19}.

%% ============================================================
\section{Accuracy Benchmarks}
\label{sec:accuracy}

\subsection{LiH potential energy surface}

We characterize the balanced coupled Hamiltonian by computing
4-electron FCI energies across the LiH potential energy surface.
Table~\ref{tab:pes_nmax2} shows the FCI energy at seven bond lengths
for $n_{\max}=2$ ($Q=30$ qubits, 15 spatial orbitals, FCI dimension
11{,}025).

\begin{table}[h]
\caption{LiH balanced coupled FCI energies at $n_{\max}=2$
  ($Q=30$, 878 Pauli terms, analytical V$_{ne}$ integrals).
  $E_{\text{exact}}$ is the non-relativistic Born-Oppenheimer
  energy~\cite{Tung2011}.}
\label{tab:pes_nmax2}
\begin{ruledtabular}
\begin{tabular}{ccccc}
$R$ (bohr) & $E_{\text{FCI}}$ (Ha) & $\Delta E$ (mHa) & Error (\%) \\
\colrule
2.0  & $-7.8246$ & 181 & 2.3 \\
2.5  & $-7.8940$ & 164 & 2.1 \\
3.0  & $-7.9290$ & 142 & 1.8 \\
3.015 & $-7.9294$ & 141 & 1.8 \\
3.5  & $-7.9273$ & 137 & 1.7 \\
4.0  & $-7.8952$ & 157 & 2.0 \\
5.0  & $-7.7716$ & 254 & 3.3 \\
\end{tabular}
\end{ruledtabular}
\end{table}

The FCI-predicted equilibrium is $R_{\text{eq}} = 3.227$~bohr
(7.0\% error vs.\ experimental 3.015~bohr). At $n_{\max}=3$
($Q=84$, 42 spatial orbitals, FCI dimension 741{,}321), the energy
at $R=3.015$~bohr improves to $-8.0545$~Ha (0.20\% error), and
the equilibrium shifts to $R_{\text{eq}} = 3.28$~bohr (8.8\% error).
Table~\ref{tab:convergence} summarizes the basis convergence.

\begin{table}[h]
\caption{Basis convergence: LiH balanced coupled at $R=3.015$~bohr.
  Exact non-relativistic energy: $-8.0706$~Ha.}
\label{tab:convergence}
\begin{ruledtabular}
\begin{tabular}{cccccc}
$n_{\max}$ & $Q$ & $E_{\text{FCI}}$ (Ha) & Error (\%) & $R_{\text{eq}}$ (bohr) & $R_{\text{eq}}$ err \\
\colrule
2 & 30  & $-7.929$ & 1.8 & 3.227 & 7.0\% \\
3 & 84  & $-8.055$ & 0.20 & 3.28 & 8.8\% \\
\end{tabular}
\end{ruledtabular}
\end{table}

The 9$\times$ improvement in energy error from $n_{\max}=2$ to
$n_{\max}=3$ demonstrates rapid basis convergence. The equilibrium
geometry drifts by $+0.053$~bohr per unit increase in $n_{\max}$,
approximately 3$\times$ smaller than the drift observed with the
PK pseudopotential~\cite{GeoVac_Paper19}. The energy error at
$n_{\max}=3$ (0.20\%) is competitive with Gaussian double-zeta
bases for single-point calculations.

\subsection{Multi-geometry FCI: solver vs.\ Hamiltonian accuracy}

The 5.3--26\% equilibrium geometry errors reported in the classical
solver literature for GeoVac~\cite{GeoVac_Paper17} arise from the
classical solver (adiabatic approximation, PK overcounting), not from
the Hamiltonian operator itself. When the balanced coupled Hamiltonian
is solved exactly via FCI:

\begin{itemize}
\item The energy error at equilibrium drops from 15\% (composed + PK)
  to 1.8\% ($n_{\max}=2$) or 0.20\% ($n_{\max}=3$).
\item The balanced coupled configuration is the \emph{only} one that
  produces a bound LiH molecule in 4-electron FCI---decoupled (no
  cross terms), composed with PK, and coupled without cross-center
  V$_{ne}$ all give unbound results~\cite{GeoVac_Paper19}.
\item The variational bound is respected at all geometries tested
  ($E_{\text{FCI}} > E_{\text{exact}}$ everywhere).
\end{itemize}

\subsection{BeH$_2$ and H$_2$O}

For BeH$_2$ ($Q=50$), the balanced coupled 4-electron FCI gives
$E=-14.146$~Ha (10.7\% error), compared to $-11.497$~Ha (27.4\%)
for composed with PK.  The balanced 1-norm (304.7~Ha) is 14\% lower
than composed (354.9~Ha), demonstrating that the PK-free architecture
is cheaper for quantum simulation.

For H$_2$O ($Q=70$), FCI is infeasible (10 electrons), but the
Hamiltonian metrics are available: 5{,}798 Pauli terms, 1{,}509~Ha
1-norm, 168 QWC groups.  The balanced 1-norm is $0.054\times$ the
composed total (which is dominated by the PK barrier at $Z_{\text{eff}}=6$),
and $4.2\times$ the composed electronic-only 1-norm.

\subsection{Applicability: single-point vs.\ dissociation}
\label{sec:applicability}

The balanced coupled Hamiltonian is designed for single-point energy
calculations at fixed molecular geometries---the use case for quantum
resource estimation, VQE benchmarks, and QPE demonstrations.  For this
primary use case, the balanced coupled Hamiltonian achieves 0.20\%
energy accuracy at $n_{\max}=3$.

The Hamiltonian is \emph{not} designed for computing dissociation curves
or reaction energies between different molecular species.  The
cross-center $V_{ne}$ multipole expansion includes a monopole ($1/R$)
component that decays slowly, causing the electronic energy to remain
unconverged even at $R = 100$~bohr.  The practical dissociation energy
computed from the PES well depth ($D_e = 0.161$~Ha at $n_{\max}=2$)
overestimates the experimental value (0.092~Ha) by 74\%.  Users
requiring dissociation energies should use Gaussian-basis methods for
the dissociation limit and GeoVac Hamiltonians for single-point
calculations at experimentally known or DFT-optimized geometries.

\subsection{Chemistry-accuracy scope boundary}
\label{subsec:scope_boundary}

The framework's PES-topology accuracy is row-conditional.  For first-row
hydrides (Z = 3--9 + H), the balanced coupled Hamiltonian binds at the
correct equilibrium geometry with documented overbinding magnitude.  For
second-row and below (Z $\ge$ 11), it does not.

\paragraph{Empirical evidence.}  The sprint sequence of
2026-06-07 (project memos
\texttt{debug/sprint\_w1e\_projection\_audit\_memo.md},
\texttt{debug/lih\_kwarg\_sweep\_log.txt}, and
\texttt{debug/sprint\_b1\_nah\_hf\_verdict\_memo.md}) systematically
tested the chemistry-accuracy scope under every available kwarg
combination and one architectural extension.  Table~\ref{tab:scope}
summarises the verdicts.

\begin{table}[h]
  \centering
  \caption{PES topology under exhaustive engineering tests at
    $n_{\max}=2$.  ``Interior min'' indicates whether the qubit-FCI or
    qubit-HF PES exhibits a binding well within R $\in$ [2.5, 6.0] bohr.}
  \label{tab:scope}
  \begin{tabular}{lcc}
    \toprule
    Configuration & LiH & NaH \\
    \midrule
    composed (PK-partitioned) & no & no \\
    balanced (frozen-core, Path A) & yes ($R_{\mathrm{eq}}=3.227$, 7\% high) & no \\
    balanced + multi-zeta valence & N/A (Z $<$ 11) & no \\
    balanced + screened cross-center & yes (unchanged) & no \\
    balanced + cross-block h1 & catastrophic (D$_e$ = 1.1 Ha) & N/A \\
    explicit-core HF (12 electrons NaH; 4 e LiH cross-check) & yes (binds) & no \\
    \bottomrule
  \end{tabular}
\end{table}

\paragraph{Operator-level diagnosis.}  Three structural factors converge
to set the scope boundary at the second-row hydride threshold:
\begin{itemize}
  \item \emph{Fixed-basis projection.}  The qubit Hamiltonian uses a
    static $Z_{\text{eff}}$ hydrogenic orbital basis at each $R$, while
    the binding physics for second-row hydrides (Na $3s$ at
    $Z_{\text{eff}} \approx 2.2$) requires R-adaptive basis content
    that is structurally absent in second-quantisation.
  \item \emph{Multipole-expansion regime.}  Track CD's cross-center
    V$_{ne}$ at $L_{\max}=2$ captures the binding-driving R-dependence
    for first-row systems (LiH $Z_{\text{eff}} \approx 1.84$), where
    the valence orbital is compact enough that fixed-basis multipole
    matrix elements track the physics.  For second-row, the more
    diffuse valence orbital outruns the multipole machinery.
  \item \emph{Frozen-core projection.}  Sprint B.1 (2026-06-07) tested
    Hartree--Fock on an explicitly un-frozen [Ne] core (12-electron
    NaH at 30 qubits, 15 spatial orbitals) and reproduced the
    monotone-overattraction pattern.  Frozen-core projection is
    \emph{not} the wall; the structural gap survives un-freezing.
\end{itemize}

\paragraph{Implication for users.}  Users running the framework for
first-row hydride chemistry can expect bound PES with quantitative
overbinding (the 0.20\% / $D_e \times 2.4$ pattern documented in
\S\ref{sec:applicability}).  Users running second-row or below should
treat the qubit Hamiltonian as a \emph{resource-estimation target}
only, not as a chemistry-accuracy Hamiltonian.  The Pauli counts,
1-norm bounds, and GH-convergence error bounds are correct in either case
(they are properties of the operator, not its physics fidelity), so
algorithm research that benchmarks against scaling behaviour transfers
cleanly across the scope boundary.  Algorithm research that wants
ground-state energies should stay in the first-row library.

\paragraph{Structural reason.}  The scope boundary has a structural
origin in the framework's NCG composition operation.  We verified
empirically on H$_2$ (2026-06-07, project memo
\texttt{debug/bratteli\_h2\_pilot\_memo.md}) that the Track CD
one-body Hamiltonian is bit-exactly (residual $9.7 \times 10^{-18}$) a
Marcolli--van Suijlekom 2014 gauge network~\cite{marcolli_vs2014} on
the molecular bond quiver, with atomic Camporesi--Higuchi Dirac data
at vertices and multipole-V$_{ne}$ bimodules on edges.  The
scope-boundary structural mechanism --- why first-row binds and
second-row does not --- is the chemistry-side instance of the
calibration-data partition documented in Paper~18, \S~IV.6.  A
forthcoming math.OA paper formalises the NCG correspondence; the
present paper takes the structural framing as a given and documents
the empirical scope.

%% ============================================================
\section{Quantum Resource Comparison}
\label{sec:resources}

\subsection{Pauli terms and measurement cost}

Table~\ref{tab:resources} compares GeoVac Hamiltonians against
Gaussian raw Jordan-Wigner baselines.

\begin{table*}[t]
\caption{Quantum resource comparison: GeoVac vs.\ Gaussian raw JW.
  GeoVac composed values use PK classical partitioning
  (electronic-only 1-norm); the GeoVac--Gaussian Pauli/QWC multipliers
  (e.g.\ $2.7\times$) use the pair-diagonal ERI rule (Sec.~\ref{sec:eri_rule})
  and re-price toward parity under the exact rule.  \textbf{Qubit-reduction
  convention:} the LiH STO-3G count (907, $Q=12$) is the \emph{raw} JW
  Hamiltonian, matched against the raw (un-tapered) GeoVac composed count
  (334, $Q=30$); the larger-basis LiH counts and all H$_2$O Gaussian counts
  apply the standard 2-qubit reduction (the $Q=20,36,12,46$ sizes are the
  reduced sizes, matching \texttt{GAUSSIAN\_LIH\_PUBLISHED}).  Under a
  uniform 2-qubit-reduced convention the headline LiH multiplier narrows
  toward parity (qubit-reduced STO-3G LiH is 276 at $Q=10$ vs.\ GeoVac
  composed 334); the raw-vs-raw and reduced-vs-reduced comparisons bracket
  the genuine advantage.  The Gaussian Pauli
  counts are the published values of Trenev et al.~\cite{Trenev2025}
  (Table~5, Appendix~B; Jordan--Wigner with 2-qubit reduction); the
  $Q^{3.9\text{--}4.3}$ Gaussian-scaling exponents are GeoVac's own log-log
  fit of those published counts.
  ``---'' indicates data not available.}
\label{tab:resources}
\begin{ruledtabular}
\begin{tabular}{llccccc}
System & Method/Basis & $Q$ & Pauli terms & 1-norm (Ha)
  & QWC groups & Advantage$^a$ \\
\colrule
\multirow{5}{*}{LiH}
 & GeoVac composed & 30 & 334 & 33.3 & 21 & --- \\
 & GeoVac balanced & 30 & 878 & 74.1 & 87 & --- \\
 & Gaussian STO-3G & 12 & 907 & 34.3 & 273 & $2.7\times$ Pauli, $13\times$ QWC \\
 & Gaussian 6-31G  & 20 & 5{,}851 & --- & --- & $17.5\times$ Pauli \\
 & Gaussian cc-pVDZ & 36 & 63{,}519 & --- & --- & $190\times$ Pauli \\
\colrule
\multirow{4}{*}{H$_2$O}
 & GeoVac composed & 70 & 778 & 361 & 21 & --- \\
 & GeoVac balanced & 70 & 5{,}798 & 1{,}509 & 168 & --- \\
 & Gaussian STO-3G & 12 & 551 & --- & --- & $0.71\times$ Pauli \\
 & Gaussian cc-pVDZ & 46 & 107{,}382 & --- & --- & $138\times$ Pauli \\
\colrule
\multirow{2}{*}{BeH$_2$}
 & GeoVac composed & 50 & 556 & 67.4 & 21 & --- \\
 & GeoVac balanced & 50 & 2{,}652 & 304.7 & 138 & --- \\
\end{tabular}
\end{ruledtabular}
\raggedright\footnotesize{$^a$Ratio of Gaussian to GeoVac composed Pauli terms.
Values $> 1$ favor GeoVac; values $< 1$ favor Gaussian.}
\end{table*}

The comparison is most favorable for LiH, where the GeoVac composed
Hamiltonian has $2.7\times$ fewer Pauli terms and $13\times$ fewer
QWC measurement groups than Gaussian STO-3G (these absolute multipliers
use the pair-diagonal ERI rule, Sec.~\ref{sec:eri_rule}, and the raw-JW
STO-3G count; they narrow toward parity under the exact rule \emph{or}
under a uniform 2-qubit-reduced convention---see the caption of
Table~\ref{tab:resources}), with a nearly identical
electronic 1-norm (33.3 vs.\ 34.3~Ha).  Against larger Gaussian bases,
the advantage grows: $190\times$ fewer Pauli terms than cc-pVDZ.

For H$_2$O, the composed Hamiltonian has \emph{more} Pauli terms than
STO-3G (778 vs.\ 551, advantage $0.71\times$) because the GeoVac
encoding uses 70 qubits (15 spatial orbitals per block $\times$ 5 blocks
$\times$ 2 spins) while STO-3G uses only 12 qubits (7 spatial orbitals
with frozen core).  The meaningful comparison is at comparable basis
quality: against cc-pVDZ (46 qubits), GeoVac wins by $138\times$.
The $O(Q^{2.2})$ vs.\ $O(Q^4)$ scaling guarantees that the GeoVac
advantage increases with basis size.

\subsection{Scaling exponents}

Table~\ref{tab:scaling} shows the within-molecule Pauli scaling exponents,
obtained from two-point fits at $n_{\max}=1$ and $n_{\max}=2$.

\begin{table}[h]
\caption{Within-molecule Pauli scaling: $N_{\text{Pauli}} \propto Q^{\alpha}$.
  GeoVac exponents from composed architecture (2-point fit,
  $n_{\max}=1,2$).  Gaussian exponent from STO-3G/6-31G/cc-pVDZ
  progression.}
\label{tab:scaling}
\begin{ruledtabular}
\begin{tabular}{lcc}
System & GeoVac $\alpha$ & Gaussian $\alpha$ \\
\colrule
LiH    & 2.24 & 3.9--4.3 \\
BeH$_2$ & 2.18 & --- \\
H$_2$O  & 2.22 & 3.9--4.3 \\
HF     & 2.21 & --- \\
NH$_3$  & 2.20 & --- \\
CH$_4$  & 2.19 & --- \\
\colrule
Mean   & 2.21 $\pm$ 0.02 & 3.9--4.3 \\
\end{tabular}
\end{ruledtabular}
\end{table}

The GeoVac composed scaling exponent of $\sim 2.2$ is nearly constant
across all six molecules, reflecting the universal Gaunt selection rule
structure.  The Gaussian scaling exponent of $\sim 4$ reflects the
dense four-index ERI tensor.  This difference means the GeoVac
advantage grows polynomially with system size: at $Q=100$, the
expected Pauli ratio is $(100)^{4.1}/(100)^{2.2} \approx 6{,}000\times$.

\subsection{Regime dependence}

The GeoVac advantage is strongest in the NISQ/VQE regime where Pauli
term count and QWC measurement groups dominate the runtime cost.
For fault-tolerant QPE, the relevant cost metric is the Pauli 1-norm
$\lambda = \sum_j |c_j|$, which determines the Trotter step
count.\footnote{%
The GH-convergence theorem on the GeoVac spectral
triple~\cite{GeoVac_Paper38} (van Suijlekom state-space Gromov--Hausdorff
distance $\Lambda(T_{n_{\max}}, T_{S^3}) \le C_3 \gamma_{n_{\max}}$ with
$\gamma_{n_{\max}} \to 0$) admits a single-particle
Lipschitz-distortion \emph{heuristic} estimate
$\|H_\infty - H_{n_{\max}}\|_{\rm op}
\lesssim N_{\rm active}\,\gamma_{n_{\max}}\,L_H$ via the standard
Childs--Su Duhamel and $N_{\rm active}$-linearity
lift~\cite{Childs2021}.  Combining this with the commutator-tightened
Trotter bound under an even error budget produces a
\emph{feasibility certificate}: the minimum $n_{\max}$ at which a
target $\epsilon$ is achievable, separate from the Trotter step
count.  At production $n_{\max} = 2$ and target $\epsilon = 10^{-3}$
the truncation budget is overshot by three to four orders of
magnitude across LiH, BeH$_2$, H$_2$O; LiH first becomes feasible at
$n_{\max} \approx 5000$.  The heuristic is uniformly looser than
naive Suzuki--Trotter at production parameters (the natural
like-for-like comparison would require $\lambda$ evaluated at the
same large $n_{\max}$, which is computationally inaccessible and is
exactly the reason the propinquity-aware route was tried).  The
heuristic is loose because (i) $L_H = \mathrm{mean}\,|c_j|$ is a
conservative single-particle Lipschitz constant, (ii) the L2
Stein--Weiss rate on the low-harmonic subspace can be much sharper
than the all-$C^\infty(S^3)$ bound, and (iii) the Duhamel plus
$N_{\rm active}$ lift admits a factor-of-2 tightening for
antisymmetric ERIs (Paper~22 Gaunt sparsity).  Sharpening each
is a multi-step research target rather than a sprint-scale fix.
Details: \texttt{debug/sprint\_trotter\_propinquity\_memo.md}.}
At matched accuracy (GeoVac balanced $n_{\max}=3$ vs.\ Gaussian
cc-pVTZ), the 1-norm comparison is less favorable due to the larger
qubit count of the GeoVac encoding (84 vs.\ 54 qubits for LiH).
However, published double-factorization or THC $\lambda$ values for
molecules at $Q < 100$ do not exist in the literature~\cite{Lee2021,
vonBurg2021}, making a direct 1-norm comparison at the QPE operating
point infeasible.

%% ============================================================
%% Tier 2 relativistic resource benchmarks — Dirac-on-S^3 drop-in (T6).
%% ============================================================
\subsection{Relativistic (spin-ful) composed resource benchmarks}
\label{subsec:paper20_tier2}

Table~\ref{tab:paper20_tier2} reports resource metrics for the
spin-ful composed encoding of Paper~14
(Dirac-on-$S^3$ Tier 2 construction). Three single-bond hydrides
spanning $Z \in \{3, 4, 20\}$, at $n_{\max} = 2$ and $n_{\max} = 3$.
Every metric is computed from the algebraic-first pipeline described
in that section: exact-rational $X_k$ angular coefficients (sympy
\texttt{wigner\_3j}), exact-rational or machine-precision $R^k$ radial
integrals (\texttt{hypergeometric\_slater.py}), closed-form
$H_{\mathrm{SO}} = -Z^4\alpha^2(\kappa+1)/[4 n^3 l(l+\tfrac{1}{2})(l+1)]$
Breit--Pauli diagonal. No numerical quadrature enters the spinor path.

\begin{table*}[h]
\caption{Tier 2 relativistic composed resource metrics (Dirac-on-$S^3$
T3 output). $N_{\mathrm{Pauli}}^{\mathrm{rel}}$: non-identity Pauli
string count in the $(\kappa, m_j)$-labeled JW encoding. $\lambda_{\mathrm{ni}}$:
non-identity Hamiltonian 1-norm. QWC: greedy qubit-wise-commuting
group count. ``scalar'' repeats the spinless benchmark at the same
molecular spec topology. Chawla ratio = $N_{\mathrm{Pauli}}^{\mathrm{rel}} / 12\,556$,
the single calibrated Chawla RaH-18q cell~\cite{Chawla2024}.}
\label{tab:paper20_tier2}
\begin{ruledtabular}
\begin{tabular}{cccccccc}
Molecule & $n_{\max}$ & $Q$ & $N_{\mathrm{Pauli}}^{\mathrm{rel}}$ & $N_{\mathrm{Pauli}}^{\mathrm{sc}}$ & $\lambda_{\mathrm{ni}}$ (Ha) & QWC & Chawla ratio \\
\hline
LiH & 2 & 30 &  1\,413 & 333    &  40.59 &  121 & 0.113$\times$ \\
LiH & 3 & 84 & 89\,226 & 7\,878 & 218.78 & 11\,865 & 7.106$\times$ \\
BeH & 2 & 30 &  1\,413 & 333    & 143.96 &  111 & 0.113$\times$ \\
BeH & 3 & 84 & 89\,226 & 7\,878 & 413.90 & 11\,843 & 7.106$\times$ \\
CaH & 2 & 20 &    942  & 222    &  18.68 &  111 & 0.075$\times$ \\
CaH & 3 & 56 & 59\,484 & 5\,252 & 125.36 & 11\,843 & 4.738$\times$ \\
SrH & 2 & 20 &    942  & 223    &  18.68 &  111 & 0.075$\times$ \\
BaH & 2 & 20 &    942  & 223    &  18.68 &  111 & 0.075$\times$ \\
\hline\hline
\multicolumn{8}{l}{\footnotesize Chawla RaH-18q~\cite{Chawla2024}: 12\,556 Pauli (rel), 2\,740 Pauli (non-rel) (47\,099 / 4\,249 two-electron integrals), $Q = 18$.} \\
\multicolumn{8}{l}{\footnotesize SrH, BaH added Sprint 3 (v2.14.0) via [Kr] and [Xe] frozen cores (\texttt{geovac/neon\_core.py}).} \\
\multicolumn{8}{l}{\footnotesize CaH/SrH/BaH at $Q=20$ bit-identical ($N_{\mathrm{Pauli}}^{\mathrm{rel}} = 942$, $\lambda_{\mathrm{ni}} = 18.68$, QWC $= 111$): isostructural invariance.} \\
\multicolumn{8}{l}{\footnotesize Sprint 4 Track TR (v2.15.0): rel Pauli counts corrected 1.76$\times$ (n=2) / 1.92$\times$ (n=3) upward via missing jj reduced-matrix-element phase fix.} \\
\multicolumn{8}{l}{\footnotesize Chawla et al.\ SI Tables S1--S3 (per-molecule matched-$Q$ data): accessibility \textbf{deferred}.} \\
\end{tabular}
\end{ruledtabular}
\end{table*}

Reading Table~\ref{tab:paper20_tier2}:

\begin{itemize}
  \item At $n_{\max}=2$, all GeoVac molecules sit at
    $0.075$--$0.11\times$ the Chawla RaH-18q Pauli count at native
    $Q \in \{20, 30\}$. Projected to matched $Q = 18$ via the
    O($Q^{2.5}$) scaling law of Paper~14 and the
    Tier 2 $\sim 4.2\times$ rel/scalar multiplier (post-TR), the advantage is
    $17\times$--$32\times$ (CaH at $Q=20$ barely shrinks to $Q=18$, setting
    the low end; LiH/BeH at $Q=30$ set the high end). This is not
    species-matched---RaH ($Z=88$) differs from LiH/BeH/CaH, and the
    per-species Chawla cells are deferred (cf.\ Paper~14).
  \item At $n_{\max}=3$, the absolute GeoVac Pauli count overtakes
    Chawla et al.\ at matched $Q$ (LiH/BeH reach 7.1$\times$ of the RaH-18q
    number at $Q = 84$; CaH sits at $4.74\times$ at $Q = 56$). This
    is the expected sparsity-exponent regime: as $l_{\max}$ grows,
    the composed architecture's per-block sparsity advantage narrows.
    The native-$Q$ sweet spot is $n_{\max}=2$.
  \item The $\lambda_{\mathrm{ni}}$ column is non-inflated by
    relativistic coupling. For LiH $n_{\max}=3$, scalar gives
    170.59 Ha and rel gives 218.78 Ha ($+28\%$); for CaH,
    96.45 Ha $\to$ 125.36 Ha ($+30\%$). The modest ($\sim 20$--$30\%$)
    increase carries no order-of-magnitude inflation.
  \item The QWC column is inflated ($\sim 15\times$ in the worst case).
    The richer $jj$-coupled Pauli structure distributes the
    Hamiltonian across more distinct Pauli words; VQE measurement
    overhead grows accordingly. Mitigable by symmetry-compressed
    or tensor-factorized measurement strategies.
\end{itemize}

\textbf{Deferred Chawla et al.\ cells.} The per-molecule Pauli counts, 1-norm,
and QWC values for BeH, MgH, CaH, SrH, BaH at $Q = 18$ are reported
in Chawla et al.~\cite{Chawla2024} Supplemental Material Tables
S1--S3. These were not extracted in the Tier 2 sprint timeframe. The
footnote row in Table~\ref{tab:paper20_tier2} marks this as deferred;
sharpening the head-to-head at matched $Q = 18$ per-molecule is
future work, noted in Paper~14.

\textbf{Matched-$Q$ feasibility (Sprint 3).} GeoVac's native-$Q$ ladder
for the composed architecture is discrete and determined by natural
geometry, not tunable: $Q = 10$ (single block at $n_{\max}=1$),
$Q = 20$ (single-bond hydride with frozen core), $Q = 30$
(explicit-core hydride), $Q = 40$+ (polyatomics or $n_{\max}=3$).
$Q = 18$ falls \emph{between} natural steps. Three truncation options
were considered in Sprint 3 Track SU and all found unphysical:
dropping a 2p orbital breaks the angular completeness of the block;
reducing $n_{\max}$ to 1 overshoots to $Q=10$; a custom 9-orbital
spec has no principled selection criterion. The correct positioning
is therefore \emph{native-$Q$ competitive resource counts}, not
matched-$Q$ per-molecule comparison. This is a structural difference
from Gaussian-basis approaches where $Q$ is tunable by choice of
basis set (STO-3G vs cc-pVDZ vs cc-pVTZ). At native $Q=20$,
CaH/SrH/BaH relativistic Hamiltonians give $942$ Pauli terms (post-TR,
Sprint 4 v2.15.0), which is $0.075\times$ Chawla et al.'s RaH-18q ($13\times$
fewer) and preserves a sparsity win under any pessimistic $Q=20 \to Q=18$
correction.

\textbf{Accuracy caveats.} The Tier 2 spin-ful encoding inherits the
composed architecture's $5$--$26\,\%$ $R_{\mathrm{eq}}$ error from
Paper~17 and the Breit--Pauli leading-order $Z\alpha^2$ approximation
from Paper~14. Chawla et al.'s DC-UCCSD is
accuracy-targeted ($0.47\%$ vs CASCI for CaH). Table~\ref{tab:paper20_tier2}
is a \emph{resource estimation} table at fixed single-point
geometries, not a spectroscopic accuracy benchmark.

\textbf{Fine-structure sanity check.} The closed-form Breit-Pauli
$H_{\mathrm{SO}}$ reproduces the He ($2^3$P), Li ($2^2$P), and Be
($2s2p\,^3$P) fine-structure splittings to correct sign and correct
order of magnitude (benchmark code
\texttt{benchmarks/fine\_structure\_check.py}). Relative errors
$-66\%, +211\%, -78\%$ against NIST~\cite{NIST_ASD} --- consistent
with the leading-order $Z\alpha^2$ approximation and simple
$Z_{\mathrm{eff}}$ screening, not spectroscopic. The 20--50\%
accuracy target suggested in the Tier 2 sprint plan is not met; full
Dirac-Coulomb radial corrections and multi-electron Breit-Pauli
corrections are Tier 3+.

%% ============================================================
\section{Available Systems}
\label{sec:systems}

The GeoVac Hamiltonian library comprises 37 systems (35 composed
molecules plus the He and H$_2$ reference atoms), spanning three rows of
the periodic table ($Z = 1$--36); Tables~\ref{tab:molecules}--%
\ref{tab:transition_metals} enumerate the molecular systems.
Table~\ref{tab:multicenter} lists the 5 multi-center molecules added
in v2.2: heteronuclear diatomics (LiF, CO), homonuclear
diatomics (N$_2$, F$_2$), and a mixed frozen-core system (NaCl).
Table~\ref{tab:transition_metals} lists all ten first-row transition metal
hydrides (v2.8), which use [Ar] frozen cores and isolated
$d$-orbital valence blocks spanning $Z = 21$--30.

\begin{table}[h]
\caption{GeoVac Hamiltonian library at $n_{\max}=2$.  First-row
  molecules use composed encoding (PK classically partitioned);
  second- and third-row molecules use balanced coupled encoding
  with frozen cores and cross-center $V_{ne}$.
  $\lambda_{\text{ni}}$ is the non-identity 1-norm (simulation-relevant).
  Isostructural pairs (rows 2--3) produce identical Pauli counts.
  The \texttt{hamiltonian()} API currently returns composed encoding
  for all rows; balanced data shown here are from Paper~19's
  \texttt{build\_balanced\_hamiltonian()}.}
\label{tab:molecules}
\begin{ruledtabular}
\begin{tabular}{llrrrr}
Molecule & Core & $Q$ & Pauli & $\lambda_{\text{ni}}$ (Ha) & QWC \\
\colrule
LiH      & He   & 30 &    334 &  26.6 & 21 \\
BeH$_2$  & He   & 50 &    556 &  52.1 & 21 \\
H$_2$O   & He   & 70 &    778 & 175.6 & 21 \\
HF       & He   & 60 &    667 & 215.1 & 21 \\
NH$_3$   & He   & 80 &    889 & 149.4 & 21 \\
CH$_4$   & He   & 90 & 1{,}000 & 129.1 & 21 \\
\colrule
NaH      & Ne   & 20 &    239 &  18.9 &  29 \\
MgH$_2$  & Ne   & 40 & 1{,}501 & 103.2 & 109 \\
HCl      & Ne   & 50 & 2{,}936 & 798.8 & 133 \\
H$_2$S   & Ne   & 60 & 4{,}119 & 810.9 & 141 \\
PH$_3$   & Ne   & 70 & 5{,}582 & 829.0 & 159 \\
SiH$_4$  & Ne   & 80 & 7{,}273 & 848.6 & 169 \\
\colrule
KH       & Ar   & 20 &    239 &  27.5 &  29 \\
CaH$_2$  & Ar   & 40 & 1{,}501 & 119.8 & 114 \\
HBr      & Ar$d$ & 50 & 2{,}936 & 807.2 & 132 \\
H$_2$Se  & Ar$d$ & 60 & 4{,}119 & 833.6 & 142 \\
AsH$_3$  & Ar$d$ & 70 & 5{,}582 & 886.2 & 159 \\
GeH$_4$  & Ar$d$ & 80 & 7{,}273 & 813.3 & 167 \\
\end{tabular}
\end{ruledtabular}
\end{table}

\begin{table}[h]
\caption{Multi-center molecules added in v2.2 at $n_{\max}=2$.
  LiF, CO, N$_2$, and F$_2$ use He frozen cores; NaCl uses mixed
  frozen cores (Na [Ne] + Cl [Ne]).}
\label{tab:multicenter}
\begin{ruledtabular}
\begin{tabular}{llcrrrr}
Molecule & Core & $N_e$ & Blocks & $Q$ & Pauli & QWC \\
\colrule
LiF      & He   & 12 & 6 &  70 &    778 & 21 \\
CO       & He   & 14 & 7 & 100 & 1{,}111 & 21 \\
N$_2$    & He   & 14 & 7 & 100 & 1{,}111 & 21 \\
F$_2$    & He   & 18 & 9 & 100 & 1{,}111 & 21 \\
NaCl     & Ne+Ne & 8 & 4 &  50 &    556 & 21 \\
\end{tabular}
\end{ruledtabular}
\end{table}

\begin{table}[h]
\caption{First-row transition metal hydrides ($Z = 21$--30) at $n_{\max}=2$
  (bond) / $n_{\max}=3$ (d-block).
  All use [Ar] frozen cores (18 electrons) with $d$-orbital valence
  blocks ($l_{\min} = 2$).  The Pauli/$Q$ ratio of~9.23
  falls below the main-group coefficient of~11.10 \emph{under the
  pair-diagonal ERI approximation} (Sec.~\ref{sec:eri_rule}); this is an
  artifact of dropping more $m$-swap integrals at higher angular momentum,
  not a sparser exact selection rule---under the exact global-$M_L$ rule the
  $d$-block ratio is \emph{higher} (30.0 vs.\ 27.9).
  Isostructural invariance holds: all ten hydrides share the same
  block topology and produce identical Pauli counts.
  Cr and Cu have anomalous configurations ($3d^5 4s^1$ and $3d^{10} 4s^1$).}
\label{tab:transition_metals}
\begin{ruledtabular}
\begin{tabular}{llcrrr}
Molecule & Config.\ & Core & $Q$ & Pauli & Pauli/$Q$ \\
\colrule
ScH & $3d^1 4s^2$ & Ar & 30 & 277 & 9.23 \\
TiH & $3d^2 4s^2$ & Ar & 30 & 277 & 9.23 \\
VH  & $3d^3 4s^2$ & Ar & 30 & 277 & 9.23 \\
CrH & $3d^5 4s^1$ & Ar & 30 & 277 & 9.23 \\
MnH & $3d^5 4s^2$ & Ar & 30 & 277 & 9.23 \\
FeH & $3d^6 4s^2$ & Ar & 30 & 277 & 9.23 \\
CoH & $3d^7 4s^2$ & Ar & 30 & 277 & 9.23 \\
NiH & $3d^8 4s^2$ & Ar & 30 & 277 & 9.23 \\
CuH & $3d^{10} 4s^1$ & Ar & 30 & 277 & 9.23 \\
ZnH & $3d^{10} 4s^2$ & Ar & 30 & 277 & 9.23 \\
\end{tabular}
\end{ruledtabular}
\end{table}

\subsection{Nuclear scaling: linearity in qubit count}

A key practical result visible in Tables~\ref{tab:molecules}
and~\ref{tab:multicenter} is that the composed Pauli term count is
\emph{exactly linear} in the qubit count $Q$ across the composed molecular library:
\begin{equation}
  N_{\text{Pauli}} = 11.10 \times Q\,.
  \label{eq:linear}
\end{equation}
This linearity arises because the composed architecture constructs
each electron-pair block independently with the same internal structure
(same angular basis, same Gaunt selection rules), so adding an atom
to a molecule adds a fixed number of blocks and a proportional number
of qubits and Pauli terms.  The practical consequence is that the
quantum measurement cost of a GeoVac Hamiltonian grows linearly with
system size---adding a CH$_3$ group to a molecule adds exactly 30
qubits and $\sim$333 Pauli terms, regardless of molecular context.
This linear nuclear scaling is a stronger result than the
$O(Q^{2.5})$ within-molecule basis scaling (Table~\ref{tab:scaling}),
which describes how cost grows when expanding the angular basis on a
\emph{fixed} molecule.

\subsection{The pair-diagonal ERI approximation: sparsity vs.\ accuracy}
\label{sec:eri_rule}

The absolute Pauli counts and multipliers in this library are realized under a
\emph{pair-diagonal} angular selection rule ($m_a = m_c$, $m_b = m_d$), a
deliberate sparsifying approximation:\ the production builders drop the
genuinely-nonzero \emph{$m$-swap} two-electron integrals that the exact
global-$M_L$ Coulomb rule ($m_a + m_b = m_c + m_d$) keeps.  Its energy effect is
sub-chemical ($\sim$\,mHa), and it is the right choice when sparsity---not energy
accuracy---is the purpose.  The approximation re-prices the absolute counts by a
constant factor per valence class---$2.51\times$ (main group), $3.25\times$
($d$-block)---and \emph{not} the $O(Q^{2.5})$ / linear scaling (a constant
prefactor leaves the slope unchanged).  Under the exact rule the LiH
small-molecule market test is near parity with STO-3G (838 vs.\ 907) and the
$d$-block coefficient is \emph{higher} (30.0), not lower; the scaling advantage
and the large-system multipliers are unaffected.  The exact rule is the one the
framework uses for precision physics; adopting it in the QC product
(``option~B'') was considered and deliberately left on the table, because the
pair-diagonal approximation is the quality quantum-simulation approximation that
serves this library's purpose.  The realized rule is pinned by
\texttt{tests/test\_paper14\_eri\_rule.py}.

Second-row molecules are accessed identically to first-row:
\begin{lstlisting}
from geovac.ecosystem_export import hamiltonian
H = hamiltonian('NaH', R=3.566)
print(H.n_terms, H.one_norm_nonidentity)
# 239, 18.91
\end{lstlisting}

\noindent The Pauli scaling exponent $Q^{2.50}$ is confirmed across both
rows.  Second-row molecules use fewer qubits than their first-row
analogues (NaH: $Q=20$ vs.\ LiH: $Q=30$) because frozen cores are
not encoded; the universal scaling exponent confirms that structural
sparsity is a property of the Gaunt selection rules, not the core
treatment.

Transition metal hydrides (Table~\ref{tab:transition_metals}) fall
\emph{below} the 11.10 coefficient: all ten first-row hydrides have
Pauli/$Q = 9.23$ under the pair-diagonal ERI approximation
(Sec.~\ref{sec:eri_rule}).  This lower ratio is an \emph{artifact} of the
pair-diagonal rule dropping more $m$-swap integrals at higher angular
momentum: under the exact global-$M_L$ rule the $d$-block ratio is
\emph{higher} (30.0 vs.\ 27.9), not lower---so $d$-orbital blocks are not
genuinely cheaper to encode than $s$/$p$ blocks at the exact rule.

\subsection{Installation and usage}

The library is available at
\url{https://github.com/jloutey-hash/geovac}
and can be installed via:
\begin{lstlisting}[language=bash]
pip install geovac-hamiltonians
\end{lstlisting}
or from source:
\begin{lstlisting}[language=bash]
git clone https://github.com/jloutey-hash/geovac
pip install -e .
\end{lstlisting}

Ecosystem compatibility: Qiskit $\geq$1.0, OpenFermion $\geq$1.6,
PennyLane $\geq$0.35. Python $\geq$3.10.

%% ============================================================
\section{Chemistry-consumer interface}
\label{sec:hybrid_pipeline}

The library exposes a chemistry-consumer interface that is
designed to interoperate with external correlation engines
(DMRG, CCSD(T), AFQMC, VQE) without requiring those engines
to know anything about GeoVac's internal spectral-triple
construction.  Three additions ship in this release.

\paragraph{GH-convergence-derived basis-truncation error bound.}
Every \texttt{GeoVacHamiltonian} returned by the public API
carries a metadata field \texttt{propinquity\_bound} (legacy field name;
the bound itself is the van Suijlekom state-space Gromov--Hausdorff
convergence estimate of Paper~38) that
evaluates the closed-form
$\Lambda \le C_3 \cdot \gamma_{n_{\max}}$ from
\cite{GeoVac_Paper38} (Paper~38, the van Suijlekom
state-space Gromov--Hausdorff convergence theorem) at the cutoff $n_{\max}$
used to build the Hamiltonian.  The L2 central-Fejér sum-rule
gives $\gamma_{n_{\max}=2,3,4} = 2.0746 / 1.6101 / 1.3223$
with asymptotic rate $4/\pi$ (the Hopf-base measure
$\mathrm{Vol}(S^2)/\pi^2$ signature of the master Mellin
engine M1, see~\cite{GeoVac_Paper18}).  This is, to our
knowledge, the first quantitative basis-truncation error
estimate exposed at a chemistry-consumer API in any framework.
The bound is system-independent (it is a property of the
spectral-triple truncation, not of the chemistry content),
so chemistry consumers reading the metadata see the same
``this energy sits within $\gamma_{n_{\max}}$ of the
framework's $n_{\max}$-truncation limit'' message regardless
of which molecule they are running.

\paragraph{FCIDUMP export.}  The \texttt{GeoVacHamiltonian}
class now exposes a \texttt{to\_fcidump(filename)} method that
writes the un-Jordan-Wigner'd $(h_1, \mathrm{eri},
e_{\rm core})$ integrals in standard FCIDUMP format
\cite{knowles1989fcidump}.  The writer is pure-Python (no
PySCF dependency added) and produces files that read back
bit-exactly: round-trip on LiH composed at $n_{\max} = 2$
gives $\max|h_1^{\rm diff}| = 0.0$ and
$\max|\mathrm{eri}^{\rm diff}| = 0.0$; a seven-system sample
across LiH, BeH$_2$, H$_2$O, NaH, KH, MgH$_2$, CaH$_2$ gives
relative error below $10^{-10}$ across all integrals.  This
exporter unblocks DMRG (via Block2 / pyscf), CCSD(T) (via
pyscf), and AFQMC (via Auxiliary-Field Monte Carlo packages)
simultaneously.

\paragraph{Openfermion-native VQE path.}  The
\texttt{GeoVacHamiltonian.to\_openfermion()} method (already
present prior to this release) is the canonical entry point
for variational quantum eigensolver (VQE) consumers.
Empirical comparison against alternative stacks is striking:
openfermion-native UCCSD on H$_2$ at $Q = 10$ qubits reaches
error $8.6 \times 10^{-13}$~mHa via L-BFGS-B from a
Hartree--Fock initialization in 165 function evaluations
($0.92$~s wall), five orders of magnitude tighter than
published STO-3G UCCSD literature.  The same H$_2$ task under
the Qiskit-Nature stack at Qiskit~2.x reaches $87$~mHa
under \texttt{efficient\_su2 + COBYLA} or fails to converge
under \texttt{uccsd}, with per-evaluation cost approximately
$3000\times$ that of the openfermion path.  Production VQE
workflows on GeoVac Hamiltonians should use the openfermion
interface and a UCCSD-style ansatz; the qiskit-nature stack
is not recommended.  Scaling to LiH at $Q = 30$ requires
sector projection beyond the qubit Hamiltonian alone (the
$2^{30}$-dimensional statevector exceeds practical memory on
commodity hardware); the four-axis Z$_2$ tapering of
\cite{GeoVac_Paper14} (Hopf $m_\ell$ + Gaunt-parity $\ell$-Z$_2$
+ atom-swap + per-sub-block particle-conservation; \texttt{tapered
= 'full'} mode of \texttt{ecosystem\_export.hamiltonian}) saves
$33\%$ on LiH ($30 \to 20$) and $31\%$ on CH$_4$ ($90 \to 62$),
materially closing the practical-scale gap for
classical-statevector simulation in the
$\sim Q = 30$--$100$ qubit regime.  The previously-reported
$\sim 12.6\%$ Pauli reduction (Hopf only) is now $\sim 30\%$ at
the production API level.

\paragraph{Honest scope.}  The FCIDUMP exporter writes the
$(h_1, \mathrm{eri}, e_{\rm core})$ tensors that
\texttt{build\_composed\_hamiltonian} and
\texttt{build\_balanced\_hamiltonian} construct.  For systems
where the composed/balanced builder produces a binding PES
under qubit/composed FCI (H$_2$, BeH$_2$, multi-center
diatomics where the multipole expansion terminates exactly),
downstream consumers receive an integral set that reproduces
the framework's qubit-level energetics.  For systems where
the composed Hamiltonian itself does not bind under qubit FCI
(LiH composed and NaH balanced at production
$n_{\max}$, observed in Sprints R3-A and R3-B,
2026-06-07), the exported integrals reproduce the same
non-binding PES; classical correlation methods consuming the
same FCIDUMP cannot close that gap because the wall is at
the projection step from continuous Level 4 to
second-quantized integrals, not at the determinant-expansion
level (cf.~\cite{GeoVac_Paper18}, \S{}IV.6).  The interface
ships honestly; consumers can use it to confirm or falsify
the projection-step localization on systems of interest.

%% ============================================================
\section{Limitations}
\label{sec:limitations}

\paragraph{Composed partition.}
The Hamiltonian is constructed within a composed partition that assigns
each electron to a specific group (core, bond, lone pair).  This is
an approximation of the full $N$-electron Hilbert space---inter-group
correlations beyond those captured by the cross-block Coulomb and
cross-center nuclear attraction integrals are neglected.  For LiH,
the full 4-electron treatment in the same angular basis gives 63.5\%
$R_{\text{eq}}$ error (angular basis limited), while the composed
partition gives 7.0--8.8\% (basis more efficient but inter-group
antisymmetry approximate)~\cite{GeoVac_Paper17}.

\paragraph{Basis completeness.}
At $n_{\max}=2$ (the production operating point for composed
Hamiltonians), single-point energies have 1.8\% error for LiH.
At $n_{\max}=3$, this improves to 0.20\%, but at the cost of
$84$ qubits and $19{,}959$ Pauli terms (balanced coupled) or
$7{,}879$ (composed).  The $n_{\max}=3$ basis approaches
chemical accuracy for single-point energies but not for
potential energy surfaces (8.8\% $R_{\text{eq}}$ error persists
due to structural drift in the composed partition).

\paragraph{Dissociation limit.}
The dissociation limit is not converged (see Sec.~\ref{sec:applicability}).

\paragraph{Classical solver accuracy.}
The classical (non-quantum) solver results reported in
Refs.~\cite{GeoVac_Paper17,GeoVac_Paper19} (5.3--26\% $R_{\text{eq}}$
error) reflect classical solver limitations, not Hamiltonian quality.
FCI on the balanced coupled Hamiltonian gives substantially better
energies (0.20\% at $n_{\max}=3$), validating the Hamiltonian for
quantum simulation use.

\paragraph{Scope.}
First-row ($Z=1$--10) and second-row ($Z=11$--18) atoms are fully
supported.  Second-row molecules use analytical [Ne] frozen cores
with Clementi--Raimondi orbital exponents; the frozen-core energy
enters only through the identity operator and does not affect
quantum measurement cost.
The full first transition series ($Z = 21$--30) is implemented as
hydrides with [Ar] frozen cores and isolated $d$-orbital valence blocks.
$d$-Orbital blocks have a lower Pauli-per-qubit ratio (9.23 vs.\ 11.10)
\emph{under the pair-diagonal ERI approximation} (Sec.~\ref{sec:eri_rule});
under the exact global-$M_L$ rule this reverses ($d$-block denser).
The atomic classifier handles all ten elements automatically,
including the anomalous configurations of Cr ($3d^5 4s^1$) and Cu ($3d^{10} 4s^1$).
Main-group third-row atoms
($Z=19$--20 and $Z=31$--36) are supported via [Ar]
and [Ar]3$d^{10}$ frozen cores~\cite{GeoVac_ScopeBoundary}.
The library currently includes 37 systems (35 composed molecules plus
the He and H$_2$ reference atoms).

\paragraph{Qubit count.}
GeoVac Hamiltonians use more qubits than Gaussian alternatives at the
same basis quality (30 vs.\ 12 for LiH at STO-3G--equivalent accuracy).
The advantage is in Pauli count and measurement cost, not qubit economy.
For near-term hardware where measurement budget (not qubit count) is
the bottleneck, this tradeoff is favorable.

%% ============================================================
\section{Conclusion}
\label{sec:conclusion}

The GeoVac Hamiltonian library provides structurally sparse qubit
Hamiltonians for molecular quantum simulation across three rows
of the periodic table (37 systems---35 composed molecules plus He and
H$_2$---$Z = 1$--36), including multi-center systems with multiple heavy
atoms and the full first transition series as hydrides (ScH through ZnH;
their lower $d$-block Pauli/$Q$ ratio is a pair-diagonal-rule artifact,
Sec.~\ref{sec:eri_rule}, not a sparser exact selection rule).  The
structural sparsity---$O(Q^{2.5})$ Pauli scaling from Gaunt selection
rules---is confirmed universal across all three rows and all molecular
topologies, independent of the core treatment (Level~3-solved or frozen).
The composed Pauli count is exactly linear in $Q$ (coefficient
$\approx 11.10$), meaning that adding atoms to a molecule increases
quantum cost linearly.
For LiH, this yields 2.7$\times$ fewer Pauli terms
and 13$\times$ fewer measurement groups than Gaussian STO-3G, growing
to 190$\times$ fewer terms against cc-pVDZ.  These \emph{absolute} multipliers
are realized under the pair-diagonal ERI approximation
(Sec.~\ref{sec:eri_rule}); under the exact global-$M_L$ rule the LiH
small-molecule count is near parity with STO-3G (838 vs.\ 907), while the
\emph{scaling} advantage (linear vs.\ $O(Q^{4})$) and the large-system
multipliers are unaffected.  For first-row hydrides
at fixed equilibrium geometries, the balanced coupled architecture
achieves 0.20\% single-point energy error at FCI; for second-row
hydrides, the qubit Hamiltonian's potential-energy surface does not
exhibit an interior minimum at $n_{\max}=2$ under any of the engineering
configurations we tested (Section~\ref{subsec:scope_boundary}).  The
framework's primary value proposition is structural sparsity with
quantitative truncation-error bounds (Paper~38 state-space GH bound, exposed as
metadata on every Hamiltonian export), not energy-accuracy parity with
CCSD(T) across the full periodic-table-row range.

Second- and third-row molecules use analytical frozen cores
([Ne], [Ar], and [Ar]3$d^{10}$), with the frozen-core energy
isolated in the identity operator.  The non-identity 1-norm---the
simulation-relevant quantity---scales comparably to first-row
molecules.  A strong isostructural invariance is observed:
frozen-core molecules with the same block topology produce
identical Pauli counts across periodic table rows (e.g.,
NaH $=$ KH $= 239$ at $Q = 20$).
The Pauli scaling exponent $Q^{2.5}$ is universal across all three
rows, confirming that the structural sparsity advantage is driven by
Gaunt selection rules rather than any specific core treatment.

Future directions include: (i)~quantum resource estimation for
fault-tolerant algorithms using GeoVac Hamiltonians, (ii)~hardware
demonstrations on trapped-ion or superconducting platforms,
(iii)~compatibility benchmarks with tensor factorization and
active-space reduction techniques, (iv)~extension to heavier main-group atoms
($Z = 37$--54) via [Kr] and [Kr]4$d^{10}$ frozen cores, and
(v)~pushing the chemistry-accuracy scope boundary into the second-row
hydride regime, which would require structural extensions beyond the
current $n_{\max}=2$ basis (multi-zeta valence orbitals,
$R$-adaptive $Z_{\text{eff}}$, bonding-orbital Pauli repulsion against
explicit cores) and is the subject of a separate engineering arc.

The library is open-source and pip-installable. All Hamiltonians
presented in this paper are reproducible from the public repository.

\begin{acknowledgments}
This work was performed using an AI-augmented research workflow
with Anthropic Claude.
\end{acknowledgments}

\bibliography{paper_20_refs}

\end{document}
